\documentclass[11pt,a4paper]{article}

% Compilation conseillée : pdflatex documentation_generale_modele_empirique.tex
\usepackage[utf8]{inputenc}
\usepackage[T1]{fontenc}
\usepackage[french]{babel}
\usepackage{lmodern}
\usepackage{microtype}
\usepackage{geometry}
\usepackage{amsmath,amssymb}
\usepackage{siunitx}
\usepackage[version=4]{mhchem}
\usepackage{booktabs}
\usepackage{array}
\usepackage{enumitem}
\usepackage{xcolor}
\usepackage{graphicx}
\usepackage{fancyhdr}
\usepackage{hyperref}

\geometry{top=2.35cm,bottom=2.35cm,left=2.45cm,right=2.45cm}
\sisetup{locale=FR,per-mode=symbol}
\DeclareSIUnit\ppm{ppm}
\DeclareSIUnit\molecule{molécule}

\definecolor{bleuCAPECL}{RGB}{17,67,104}
\definecolor{bleupale}{RGB}{234,243,250}
\definecolor{orangealerte}{RGB}{180,92,20}
\hypersetup{
  colorlinks=true,
  linkcolor=bleuCAPECL,
  urlcolor=blue!60!black,
  pdftitle={Documentation générale du modèle empirique d'émissivité atmosphérique}
}

\pagestyle{fancy}
\fancyhf{}
\lhead{CAPECL -- Projet climat 2026}
\rhead{Modèle empirique}
\cfoot{\thepage}
\setlength{\headheight}{14pt}

\newcommand{\code}[1]{\texttt{#1}}
\newcommand{\important}[1]{%
  \par\medskip
  \noindent\fcolorbox{bleuCAPECL}{bleupale}{%
    \parbox{0.94\linewidth}{#1}}%
  \par\medskip
}

\title{\vspace{2.2cm}
  {\color{bleuCAPECL}\LARGE\bfseries Modèle empirique d'émissivité\\[0.25em]
  atmosphérique}\par\vspace{0.8cm}
  {\large Documentation générale du modèle}}
\author{Projet climat 2026\\Groupe C -- Geckoquins}
\date{Juin 2026}

\begin{document}

\maketitle
\thispagestyle{empty}
\vfill
\begin{center}
\begin{minipage}{0.78\textwidth}
\small
Cette documentation complète les programmes du dossier
\code{modèle1\_empirique\_emisivité}. Elle décrit les objets physiques,
les hypothèses et les paramètres utilisés pour représenter l'atmosphère.
Elle n'est pas une explication ligne par ligne du code.
\end{minipage}
\end{center}
\vfill
\newpage

\tableofcontents
\newpage

\section{Question physique et objectif}

La surface terrestre émet un rayonnement thermique principalement dans
l'infrarouge. Certaines longueurs d'onde sont absorbées par la vapeur d'eau,
le dioxyde de carbone, le méthane et l'ozone. Le modèle cherche à estimer,
pour une tranche d'air donnée, la fraction du rayonnement absorbée à chaque
longueur d'onde.

Le contexte général est celui du
\href{https://science.nasa.gov/earth/earth-observatory/climate-and-earths-energy-budget/}{bilan
énergétique de la Terre présenté par la NASA}. Le présent modèle n'en
reproduit cependant qu'un maillon : l'absorption infrarouge locale par les
gaz. Il ne calcule ni le flux infrarouge sortant au sommet de l'atmosphère,
ni l'évolution de la température terrestre.

La chaîne de modélisation est la suivante :

\begin{equation*}
\boxed{P(z),T(z),x_g(z)}
\longrightarrow
\boxed{n_g(z)}
\longrightarrow
\boxed{\sigma_g(\lambda)}
\longrightarrow
\boxed{\tau_g(\lambda,z)}
\longrightarrow
\boxed{\alpha(\lambda,z)}.
\end{equation*}

Le modèle empirique sert avant tout à rendre cette chaîne transparente et à
tester les ordres de grandeur avant l'emploi du modèle spectroscopique
HITRAN.

\section{Système modélisé}

L'atmosphère est découpée en couches horizontales homogènes. Une couche est
repérée par une altitude $z$ et possède une épaisseur $\Delta z$. Dans la
table produite, $\Delta z=\SI{1}{\kilo\metre}$ et $z$ varie de 0 à
\SI{100}{\kilo\metre} par pas de \SI{1}{\kilo\metre}.

Dans chaque couche, on considère constants :

\begin{itemize}[itemsep=2pt]
  \item la pression $P$ et la température $T$ ;
  \item les fractions molaires $x_g$ des quatre gaz ;
  \item la section efficace $\sigma_g$ à la longueur d'onde étudiée ;
  \item la direction de propagation, verticale, sur la distance $\Delta z$.
\end{itemize}

Le calcul est monochromatique : chaque longueur d'onde est traitée
indépendamment. La diffusion, les nuages, les aérosols, la réfraction, la
courbure terrestre et les échanges radiatifs entre couches ne sont pas
représentés.

\section{Profil vertical de l'atmosphère}

\subsection{Température et pression}

La température est affine par morceaux, avec des changements de gradient aux
altitudes 11, 20, 32, 50, 71 et \SI{85}{\kilo\metre}. Cette structure reprend
l'esprit de la
\href{https://www.ngdc.noaa.gov/stp/space-weather/online-publications/miscellaneous/us-standard-atmosphere-1976/us-standard-atmosphere_st76-1562_noaa.pdf}{U.S.
Standard Atmosphere 1976}, sans lui être identique. La zone de stratopause est
simplifiée et la croissance imposée au-dessus de \SI{85}{\kilo\metre} est une
paramétrisation propre au projet.

Dans une zone commençant à $z_0$, de température $T_0$, de pression $P_0$ et
de gradient thermique constant $\gamma=\mathrm dT/\mathrm dz$, la température
est

\begin{equation}
  T(z)=T_0+\gamma(z-z_0).
\end{equation}

La pression découle de l'équilibre hydrostatique et de la loi des gaz
parfaits. Pour $\gamma\neq0$,

\begin{equation}
  P(z)=P_0\left(\frac{T(z)}{T_0}\right)^{-gM/(R\gamma)},
\end{equation}

et, pour une zone isotherme,

\begin{equation}
  P(z)=P_0\exp\!\left[-\frac{gM(z-z_0)}{RT_0}\right].
\end{equation}

Les valeurs retenues sont $g=\SI{9.80665}{\metre\per\second\squared}$,
$M=\SI{0.0289644}{\kilogram\per\mole}$,
$R=\SI{8.31445}{\joule\per\mole\per\kelvin}$ et
$P(0)=\SI{101325}{\pascal}$.

\subsection{Composition gazeuse}

Les concentrations sont des fractions molaires exprimées en parties par
million. Pour le gaz $g$,

\begin{equation}
  x_g(z)=C_g(z)\times10^{-6},
\end{equation}

où $C_g$ est la concentration en ppm. Les profils utilisés sont résumés dans
le tableau~\ref{tab:profils_gaz}.

\begin{table}[ht]
\centering
\small
\begin{tabular}{>{\bfseries}l p{0.70\linewidth}}
\toprule
Gaz & Paramétrisation utilisée \\
\midrule
\ce{CO2} & \SI{425}{\ppm} jusqu'à \SI{85}{\kilo\metre}, puis décroissance
exponentielle de hauteur caractéristique \SI{10}{\kilo\metre}. \\
\ce{CH4} & \SI{1.9}{\ppm} jusqu'à \SI{11}{\kilo\metre}, puis décroissance
exponentielle de hauteur caractéristique \SI{25}{\kilo\metre}. \\
\ce{O3} & Fond de \SI{0.1}{\ppm}; entre 15 et \SI{35}{\kilo\metre}, pic
gaussien centré à \SI{25}{\kilo\metre} et de maximum \SI{8}{\ppm}. \\
\ce{H2O} & Humidité relative décroissante dans la troposphère; plancher de
\SI{4}{\ppm} dans la haute atmosphère. \\
\bottomrule
\end{tabular}
\caption{Profils de mélange idéalisés employés dans les deux modèles.}
\label{tab:profils_gaz}
\end{table}

Les valeurs de \ce{CO2} et de \ce{CH4} sont des valeurs de référence
arrondies, cohérentes avec l'ordre de grandeur des mesures du
\href{https://gml.noaa.gov/ccgg/}{Global Greenhouse Gas Reference Network de
la NOAA}. Elles ne constituent pas une série temporelle observée.

Pour la vapeur d'eau, la pression de vapeur saturante est approchée par une
forme de Magnus :

\begin{equation}
  e_s(T_C)=611.2\exp\!\left(\frac{17.627\,T_C}{T_C+243.04}\right),
\end{equation}

avec $T_C$ en degrés Celsius et $e_s$ en pascals. L'humidité relative est
fixée à $0.77(1-z/12)$ dans la troposphère. Cette loi d'humidité, comme les
profils de méthane et d'ozone, est un choix interne de modélisation et non un
profil météorologique mesuré.

\section{De la pression au nombre de molécules}

La loi des gaz parfaits, écrite à l'échelle moléculaire, fournit la densité
totale de l'air :

\begin{equation}
  n_{\mathrm{air}}(z)=\frac{P(z)}{k_{\mathrm B}T(z)},
  \qquad
  k_{\mathrm B}=\SI{1.380649e-23}{\joule\per\kelvin}.
\end{equation}

La valeur de $k_{\mathrm B}$ est exacte dans le Système international; voir la
\href{https://www.nist.gov/si-redefinition/kelvin/kelvin-boltzmann-constant}{présentation
du NIST sur la constante de Boltzmann}. La densité du gaz $g$ vaut alors

\begin{equation}
  n_g(z)=x_g(z)n_{\mathrm{air}}(z).
\end{equation}

$n_{\mathrm{air}}$ et $n_g$ sont exprimées en molécules par mètre cube. Cette
étape fait diminuer l'absorption avec l'altitude même lorsqu'une fraction
molaire reste constante, car la densité totale de l'air suit la pression.

\section{Représentation empirique des bandes d'absorption}

\subsection{Forme choisie}

Une bande moléculaire réelle contient de nombreuses raies. Le modèle
empirique remplace chaque grande bande par un pic exponentiel symétrique en
valeur absolue :

\begin{equation}
  \sigma_j(\lambda)=\sigma_{\max,j}
  10^{-a_j|\lambda-\lambda_{0,j}|/\lambda_{0,j}}.
\end{equation}

Pour un gaz, les pics sont additionnés :

\begin{equation}
  \sigma_g(\lambda)=\sum_j\sigma_j(\lambda).
\end{equation}

Le paramètre $a_j$ est déduit de deux bornes choisies autour du maximum :

\begin{equation}
  a_j=\frac12\left(
  \frac{1}{|\lambda_{g,j}-\lambda_{0,j}|/\lambda_{0,j}}+
  \frac{1}{|\lambda_{d,j}-\lambda_{0,j}|/\lambda_{0,j}}
  \right).
\end{equation}

La forme conserve la position et l'ordre de grandeur des grandes bandes,
mais ne reproduit ni les raies individuelles ni leur élargissement avec la
pression et la température.

\subsection{Bandes retenues}

\begin{table}[ht]
\centering
\small
\begin{tabular}{>{\bfseries}l l l}
\toprule
Gaz & Centres des pics (\si{\micro\metre})
    & Maximums saisis (\si{\centi\metre\squared\per\molecule}) \\
\midrule
\ce{CO2} & 4.2; 15.0 & $13\times10^{-18}$; $4\times10^{-18}$ \\
\ce{CH4} & 3.3; 7.7 & $1.8\times10^{-18}$; $7.7\times10^{-19}$ \\
\ce{H2O} & 2.7; 6.3; 17.4 & $723\times10^{-21}$;
$916\times10^{-21}$; $31\times10^{-21}$ \\
\ce{O3} & 8.9; 9.6; 14.3 & $2.1\times10^{-21}$;
$42\times10^{-21}$; $1.8\times10^{-21}$ \\
\bottomrule
\end{tabular}
\caption{Paramètres spectraux actuellement inscrits dans le modèle empirique.}
\label{tab:bandes}
\end{table}

Les sections saisies en \si{\centi\metre\squared\per\molecule} sont
converties par

\begin{equation}
  \SI{1}{\centi\metre\squared}=10^{-4}\,\si{\metre\squared}.
\end{equation}

\important{\textbf{Traçabilité des paramètres.} Aucune référence primaire
associée aux maximums et aux largeurs du tableau~\ref{tab:bandes} n'est
enregistrée dans le dossier actuel. Ces nombres doivent donc être considérés
comme des paramètres de calibration du prototype, et non comme des données
spectroscopiques de référence. Leur comparaison au modèle HITRAN est
indispensable avant toute interprétation quantitative.}

\section{Épaisseur optique et absorptivité}

Pour une couche homogène d'épaisseur $\Delta z$, l'épaisseur optique due au
gaz $g$ est

\begin{equation}
  \tau_g(\lambda,z)=\sigma_g(\lambda)n_g(z)\Delta z.
\end{equation}

Les contributions des absorbeurs indépendants s'additionnent :

\begin{equation}
  \tau_{\mathrm{tot}}(\lambda,z)=
  \tau_{\ce{CO2}}+\tau_{\ce{H2O}}+\tau_{\ce{CH4}}+\tau_{\ce{O3}}.
\end{equation}

La transmission monochromatique suit la loi de Beer--Lambert,

\begin{equation}
  \mathcal T(\lambda,z)=\exp[-\tau_{\mathrm{tot}}(\lambda,z)],
\end{equation}

et, en l'absence de diffusion et de réflexion, l'absorptivité vaut

\begin{equation}
  \boxed{\alpha(\lambda,z)=1-\exp[-\tau_{\mathrm{tot}}(\lambda,z)]}.
\end{equation}

Les programmes nomment cette grandeur «~émissivité~». L'identification
$\varepsilon_\lambda=\alpha_\lambda$ repose sur la loi de Kirchhoff à
l'équilibre thermodynamique local. Elle ne signifie pas que le modèle calcule
un flux émis : il faudrait pour cela la loi de Planck et la propagation du
rayonnement à travers toutes les couches. Les notions d'épaisseur optique et
de densité de colonne sont détaillées dans la page
\href{https://hitran.org/docs/definitions-and-units/}{Definitions and units de
HITRANonline}.

Deux limites simples servent de contrôle :

\begin{align}
  \tau_{\mathrm{tot}}\ll1 &\quad\Rightarrow\quad
  \alpha\simeq\tau_{\mathrm{tot}},\\
  \tau_{\mathrm{tot}}\gg1 &\quad\Rightarrow\quad
  \alpha\longrightarrow1.
\end{align}

\section{Domaine et lecture du tableau final}

La table \code{table\_emissivite\_0\_20um.csv} contient actuellement
\num{20301} lignes de données : 101 altitudes et 201 longueurs d'onde. Chaque
ligne enregistre les conditions thermodynamiques, les concentrations, les
sections efficaces, les épaisseurs optiques partielles et totale, puis
l'absorptivité appelée \code{emissivite}.

\important{Le point $\lambda=\SI{0}{\micro\metre}$ n'a pas de sens physique
et doit être supprimé lors d'une prochaine génération. De plus, les queues
exponentielles ne deviennent jamais strictement nulles : le modèle attribue
donc une faible absorption résiduelle en dehors des bandes choisies.}

Chaque ligne décrit une couche isolée. Additionner directement les
«~émissivités~» des altitudes est incorrect. Pour obtenir la transmission
d'une colonne, il faut additionner les épaisseurs optiques ou multiplier les
transmissions. Pour obtenir un rayonnement sortant, il faut en plus ajouter
l'émission propre de chaque couche et son atténuation par les couches
supérieures.

\section{Limites et usage raisonnable}

Le modèle est adapté à l'étude des dépendances et à la vérification de la
chaîne de calcul. Il ne doit pas être utilisé seul pour une prévision
climatique ou un flux radiatif absolu. Ses principales limites sont :

\begin{itemize}[itemsep=2pt]
  \item des bandes lissées et calibrées manuellement ;
  \item des sections efficaces indépendantes de $P$ et $T$ ;
  \item des profils verticaux idéalisés ;
  \item une extrapolation très simplifiée au-dessus de \SI{85}{\kilo\metre} ;
  \item des couches homogènes et un trajet uniquement vertical ;
  \item l'absence de nuages, d'aérosols, de diffusion et du continuum de
        vapeur d'eau ;
  \item l'absence de fonction de Planck, de flux montant ou descendant et de
        bilan énergétique.
\end{itemize}

La suite logique est le modèle HITRAN : il conserve la construction de
l'épaisseur optique, mais remplace les pics empiriques par des spectres issus
de paramètres de transitions moléculaires.

\section{Correspondance avec les fichiers du projet}

Cette section sert seulement de carte d'accès. Les scripts intermédiaires sont
des étapes de validation et non quatre modèles physiques indépendants.

\begin{table}[ht]
\centering
\small
\begin{tabular}{p{0.43\linewidth}p{0.49\linewidth}}
\toprule
Fichier & Objet scientifique associé \\
\midrule
\code{code1\_emissivité\_une\_couche.py} & Opacité d'une couche pour un gaz. \\
\code{code2\_emissivité\_n\_couche.py} & Addition des contributions de
plusieurs gaz dans une même couche. \\
\code{code3\_emissivite\_avec\_concentrations.py} & Variation de la
composition avec l'altitude. \\
\code{calcul\_section\_efficace\_CO2\_CH4\_H2O.py} & Enveloppes empiriques
des bandes d'absorption. \\
\code{code4\_emissivite\_avec\_section\_eff.py} & Réunion du profil vertical,
des spectres et de Beer--Lambert. \\
\code{generer\_table\_emissivite.py} & Échantillonnage du modèle dans le CSV. \\
\bottomrule
\end{tabular}
\end{table}

Malgré son nom, \code{code2\_emissivité\_n\_couche.py} traite plusieurs gaz
dans une couche; il ne propage pas encore le rayonnement à travers $n$
couches atmosphériques.

\section{Sources et ressources}

\begin{itemize}[leftmargin=*,itemsep=3pt]
  \item \href{https://science.nasa.gov/earth/earth-observatory/climate-and-earths-energy-budget/}{NASA Earth Observatory -- Climate and Earth's Energy Budget} ;
  \item \href{https://www.ngdc.noaa.gov/stp/space-weather/online-publications/miscellaneous/us-standard-atmosphere-1976/us-standard-atmosphere_st76-1562_noaa.pdf}{NASA/NOAA -- U.S. Standard Atmosphere 1976} ;
  \item \href{https://gml.noaa.gov/ccgg/}{NOAA Global Monitoring Laboratory -- Carbon Cycle Greenhouse Gases} ;
  \item \href{https://www.nist.gov/si-redefinition/kelvin/kelvin-boltzmann-constant}{NIST -- Boltzmann constant and the kelvin} ;
  \item O.~A. Alduchov et R.~E. Eskridge,
  \href{https://doi.org/10.1175/1520-0450(1996)035\%3C0601:imfaos\%3E2.0.co;2}{Improved Magnus Form Approximation of Saturation Vapor Pressure}, 1996 ;
  \item \href{https://hitran.org/docs/definitions-and-units/}{HITRANonline -- Definitions and units}.
\end{itemize}

\end{document}
